\documentclass[12pt,a4paper,figuresright]{book}

\usepackage{amsmath,amssymb}
\usepackage{tabularx,graphicx,url,xcolor,rotating,multicol,epsfig,colortbl}

\setlength{\textheight}{25.2cm}
\setlength{\textwidth}{16.5cm} %\setlength{\textwidth}{18.2cm}
\setlength{\voffset}{-1.6cm}
\setlength{\hoffset}{-0.3cm} %\setlength{\hoffset}{-1.2cm}
\setlength{\evensidemargin}{-0.3cm} 
\setlength{\oddsidemargin}{0.3cm}
\setlength{\parindent}{0cm} 
\setlength{\parskip}{0.3cm}

% -- adding a talk
% -- adding a talk
\newenvironment{talk}[9]% [1] talk title
                         % [2] speaker name, [3] affiliations, [4] email,
                         % [5] coauthors, [6] special session
                         % [7] time slot
                         % [8] talk id, [9] session id or photo
 {%\needspace{6\baselineskip}%
  \vskip 0pt\nopagebreak%
%   \colorbox{gray!20!white}{\makebox[0.99\textwidth][r]{}}\nopagebreak%
%   \ifthenelse{\equal{#11}{photo}}{%
%                     \\\\\colorbox{gray!20!white}{\makebox{\includegraphics[width=3cm]{#10}}}\nopagebreak}{}%
 \vskip 0pt\nopagebreak%
%  \label{#10}%
  \textbf{#1}\vspace{3mm}\\\nopagebreak% 
  \textit{#2}\\\nopagebreak%
  #3\\\nopagebreak%
  \url{#4}\vspace{3mm}\\\nopagebreak%
  \textit{#5}\\\nopagebreak%
  #6\\\nopagebreak%
  \url{#7}\vspace{3mm}\\\nopagebreak%
  \ifthenelse{\equal{#8}{}}{}{Coauthor(s): #8\vspace{3mm}\\\nopagebreak}%
  \ifthenelse{\equal{#9}{}}{}{Special session: #9\quad \vspace{3mm}\\\nopagebreak}%
 }
 {\vspace{1cm}\nopagebreak}%
 
\pagestyle{empty}

% ------------------------------------------------------------------------
% Document begins here
% ------------------------------------------------------------------------
\begin{document}
	
\begin{talk}
  {Investigating general $L2$ discrepancies with Message-Passing Monte Carlo}% [1] talk title
  {Ally Kwan}% [2] speaker name
  {Foothill College}% [3] affiliations
  {ally.kwan@yale.edu}% [4] email
  {Lijia Lin}% [5] speaker name 2
  {Johns Hopkins University}% [6] affiliation 2
  {llin72@jh.edu}% [7] email 2
  {Nathan Kirk}% [8] coauthors
  {}% [9] special session. Leave this field empty for contributed talks. 
			
Finding low-discrepancy (LD) points are essential to facilitate the use of quasi-Monte Carlo methods, which achieve faster convergence rate compared to traditional Monte Carlo methods. The Message-Passing Monte Carlo (MPMC)\textsuperscript {1} method utilizes graph neural networks to generate LD point sets, where, in the original model, the learning objective is set to the classical $L2$ star discrepancy. In this talk, we employ MPMC as an optimization tool to investigate the properties of LD point sets optimized with respect to other $L2$ discrepancy measures known in the literature and show evidence that the star discrepancy should not always be the standard measure to use in optimization pipelines. Further, we demonstrate the new MPMC module in the QMCPy Python package that is available for use with functionality such as choosing the discrepancy learning objective, specifying weights for high-dimensional problems, randomizations, and pretrained points for particular $N$ and $d$ combinations to facilitate use of the MPMC method.


\begin{enumerate}
	\item[{[1]}] Rusch, T. Konstantin and Kirk, Nathan and Bronstein, Michael M. and Lemieux, Christiane and Rus, Daniela (2024). {\it Message-Passing Monte Carlo: Generating low-discrepancy point sets via graph neural networks}. Proceedings of the National Academy of Sciences (PNAS).

\end{enumerate}

\end{talk}

\end{document}

